\section{What Survives: Coherence vs.\ Similarity}
\label{sec:coherence_survives}

The matched-similarity test rejects the strong CCS interpretation, but it
should not be overread as evidence that context-set structure is useless. We
compare two kinds of uninformative context: random off-topic noise and
same-topic passages that preserve some contextual coherence while omitting the
answer. If coherence were only generic retrieval quality under another name,
the slopes should look similar.

\begin{table}[H]
\centering
\small
\caption{Noise slope response. Coherence-preserving uninformative noise hurts
less than random off-topic noise at matched rate.}
\label{tab:noise}
\begin{tabular}{lrrr}
\toprule
Condition & Faith slope & Similarity slope & Drop at full noise \\
\midrule
Random noise & -0.069 & -0.481 & 0.052 \\
Coherent uninformative noise & -0.043 & -0.113 & 0.042 \\
\bottomrule
\end{tabular}
\end{table}

This supports a narrower claim. Coherence carries diagnostic signal beyond
mean query similarity, but it does not identify support by itself. A coherent
context set can still omit the answer, and a low-coherence set can still
contain the necessary evidence.
